% ════════════════════════════════════════════════════════════════════════════════
\section{Multi-Factor Authentication}
\label{sec:mfa}

We implemented three authentication factors in order of increasing
robustness: a password challenge-response scheme, TOTP, and FIDO2/WebAuthn.
Which factors are required for a given login depends on the RBA score
(Section~\ref{sec:pam}).

\subsection{Helpdesk Enrolment}
\label{sec:helpdesk}
Credential management is handled by a separate Flask app,
\texttt{idp/helpdesk/app.py}, bound to \texttt{127.0.0.1:8081} with
\texttt{network\_mode: host}.
Being bound to the loopback interface, it is physically unreachable from
outside the gateway machine. Its access boundary is enforced by the
network, not by authentication logic.

The helpdesk shares the same SQLite database as the IdP through a
bind-mount, so credentials registered here become immediately usable at
login. It exposes endpoints for all three factors
(\texttt{/user/<u>/password}, \texttt{/totp}, \texttt{/fido2/start},
\texttt{/fido2/finish}) plus credential listing and removal.

\subsection{Password (Challenge-Response)}
\label{sec:mfa-pwd}
The raw password never leaves the browser.

\textbf{Enrolment:} 

The helpdesk derives a 32-byte verifier with
scrypt~\cite{rfc7914,percival2009stronger}:
\[
v = \text{scrypt}(\mathit{pw},\; s,\; N=2^{14},\; r=8,\; p=1,\; \text{dkLen}=32)
\]
where $\mathit{pw}$ is the user's password and $s$ is a 16-byte random
salt; only $(s, v)$ are stored in the \texttt{passwords} table.

\textbf{Login:} 

When the PAM module requests the password factor, the
browser fetches \texttt{/oauth/mfa/<sid>/password/challenge}, which returns
the salt, a fresh 32-byte nonce, and the scrypt parameters.
A pure-JavaScript scrypt implementation
(\texttt{idp/static/js/scrypt.js}) recomputes $v'$ in the browser and
submits
\[
\text{proof} = \text{HMAC-SHA256}(v',\; \text{nonce})
\]
The server compares against \texttt{HMAC-SHA256(v, nonce)} with
\texttt{hmac.compare\_digest}.
The nonce is burned atomically (\texttt{BEGIN IMMEDIATE}) before the
comparison, so a replay finds the nonce already consumed.
Fig.~\ref{fig:mfa-password} shows the browser form during this step.

\begin{figure}[H]
  \centering
  % HOW TO TAKE: capture the browser password challenge form during a
  % VPN login. The form shows a password input; the page title should
  % indicate it is the "Password" step of the MFA flow.
  \imgplaceholder[3.5cm]{Browser: password challenge-response form during VPN login}
  \caption{Password step: scrypt and the HMAC proof run client-side.}
  \label{fig:mfa-password}
\end{figure}

\subsection{TOTP}
\label{sec:mfa-totp}
RFC 6238~\cite{rfc6238}, six digits, 30-second step, $\pm 1$ window, via
\texttt{pyotp.TOTP.verify(code, valid\_window=1)}.

Enrolment happens at the helpdesk:
\texttt{POST /user/<u>/totp} generates a fresh secret and shows it as a QR
code (Fig.~\ref{fig:totp-enrol}). Confirmation requires the user to submit
a valid code, after which the row in \texttt{totp\_secrets} flips to
\texttt{confirmed=1}.

The secret is stored encrypted at rest under AES-256-GCM
(\texttt{common/db.py:\_encrypt}); see Section~\ref{sec:security}.
At login the browser presents a 6-digit code entry form
(Fig.~\ref{fig:mfa-totp}).

\begin{figure}[H]
  \centering
  \begin{minipage}[b]{0.47\columnwidth}
    \imgplaceholder[3cm]{Helpdesk: TOTP QR code displayed during enrolment}
    \caption{TOTP QR code during helpdesk enrolment.}
    \label{fig:totp-enrol}
  \end{minipage}
  \hfill
  \begin{minipage}[b]{0.47\columnwidth}
    % HOW TO TAKE: capture the 6-digit TOTP code input form that appears
    % in the browser during a VPN login (after the password step).
    \imgplaceholder[3cm]{Browser: TOTP 6-digit code entry during VPN login}
    \caption{TOTP step during VPN authentication.}
    \label{fig:mfa-totp}
  \end{minipage}
\end{figure}

\subsection{FIDO2 / WebAuthn}
\label{sec:mfa-fido}
We use the \texttt{webauthn} (v$\geq$2) Python library on the server and the
browser's native \texttt{navigator.credentials} API on the client.

Enrolment uses the helpdesk's \texttt{/fido2/start} and
\texttt{/fido2/finish} endpoints; the browser's WebAuthn dialog asks
the user to select the authenticator device (Fig.~\ref{fig:fido2-reg}),
and on success the credential appears in the user's panel
(Fig.~\ref{fig:fido2-done}).

Both registration and authentication set
\texttt{UserVerificationRequirement.REQUIRED}, which forces the
authenticator to validate a local PIN, biometric, or password before
releasing the credential's private key.
A physically stolen, unlocked security key therefore cannot satisfy this
factor (the assertion will be missing the \texttt{uv} flag, and the server
rejects it).

At login the IdP issues a challenge via
\texttt{/oauth/mfa/<sid>/fido2/start} (bound to the Flask session), the
browser-side glue (\texttt{idp/static/js/webauthn.js}) calls
\texttt{navigator.credentials.get()}, and the assertion is verified with
\texttt{verify\_authentication\_response(\dots,
require\_user\_verification=True)}.
The credential's sign counter is updated on every successful assertion to
detect cloned authenticators.

\begin{figure}[H]
  \centering
  \begin{minipage}[b]{0.47\columnwidth}
    \imgplaceholder[3cm]{Browser: ``Choose where to save your passkey'' dialog showing USB security key option}
    \caption{Browser authenticator selection during FIDO2 registration.}
    \label{fig:fido2-reg}
  \end{minipage}
  \hfill
  \begin{minipage}[b]{0.47\columnwidth}
    \imgplaceholder[3cm]{Helpdesk: FIDO2 credential listed after successful registration}
    \caption{Helpdesk: FIDO2 credential listed after registration.}
    \label{fig:fido2-done}
  \end{minipage}
\end{figure}
